\documentclass[addpoints]{exam}
% \documentclass[addpoints,12pt]{exam}

\usepackage[slovene]{babel}
\usepackage{amsfonts}
\usepackage{amssymb}
\usepackage{amsmath}
\usepackage{float}
\usepackage{graphicx}
\usepackage{enumitem}
\usepackage{color}
\usepackage{eurosym}


%Komentar naslednje vrstice glede na verzijo testa -- rešitve ali prostor za odgovore:
% \printanswers

% \linespread{2}


\pointsinrightmargin
\marginbonuspointname{}


\renewcommand{\solutiontitle}{\noindent\textbf{Rešitve in točkovanje:}\par\noindent}

\colorgrids
\definecolor{GridColor}{gray}{0.7}
\setlength{\gridsize}{7mm}

\extrawidth{5mm}
\setlength{\rightpointsmargin}{1.5cm}

\begin{document}

\runningfooter{}{}{}

\begin{center}
    \fbox{\fbox{\parbox{15cm}{\centering
    Drugo pisno ocenjevanje znanja \\ 1. letnik -- test A \\ 4. december 2025 \\ (Potence in izrazi; Deljivost)}}}
\end{center}

\vspace{0.1in}
\makebox[\textwidth]{Ime in priimek, oddelek:\enspace\hrulefill}


\begin{center}
    \chqword{Naloga}
    \chpword{Možne točke}
    \chbpword{Dodatne točke}
    \chsword{Dosežene točke}
    \chtword{Skupaj}  
    \combinedpointtable[h][questions]
    % \combinedgradetable[h][questions]

\end{center}

\vspace{0.1in}


\begin{questions}

    % 1. vprašanje
    
    \question
    Izračunajte.
\begin{parts}
    \part[4]
    $\left(\left(-2\right)^3+\left(-5\right)^2\right)\cdot\left(-1\right)^{2025}-\left(\left(-3\right)^2\cdot\left(-2\right)-\left(-4\right)^2\right)-1^{2026}$

    \begin{solution}[4.5cm]
    $$\begin{aligned}
    \left(\left(-2\right)^3+\left(-5\right)^2\right)\cdot\left(-1\right)^{2025}-\left(\left(-3\right)^2\cdot\left(-2\right)-\left(-4\right)^2\right)-1^{2026}&=(-8+25)(-1)-(9(-2)-16)-1= \\ &=-17+34-1= \\ &=-18
    \end{aligned}$$
        Za pravilno uporabo pravila potenciranja z negativnim predznakom $1$ točka, pravilno potenciranje $1$ točka, pravilno množenje $1$ točka, končen rezultat $1$ točka.
    \end{solution}



    \part[3]
    $5\cdot 2^{2n+1}-2\cdot 4^{n+1}-2\cdot 2^{2n+2}$

    \begin{solution}[4.5cm]
    $$5\cdot 2^{2n+1}-2\cdot 4^{n+1}-2\cdot 2^{2n+2}=5\cdot 2^{2n+1}-2^{2n+3}-2^{2n+3}=2^{2n+1}(5-2^2-2^2)=-3\cdot 2^{2n+1}$$
        Za pravilen zapis potenc števila $2$ $1$ točka, pravilno izpostavljanje skupnega faktorja $1$ točka, končen rezultat $1$ točka.
    \end{solution}

        
    \end{parts}


    % 2. vprašanje
    \question[5]
    Skrčite izraz in rezultat razstavite: $(7a-2b)^2-(5a+2b)(5a-2b)-6(4a^2-b^2)$.

    \begin{solution}[5cm]
        \begin{align*}
            (7a-2b)^2-(5a+2b)(5a-2b)-6(4a^2-b^2)&=49a^2-28ab+4b^2-25a^2+4b^2-24a^2+6b^2=\\
                &=14b^2-28ab= \\
                &=14b(b-2a)
        \end{align*}
        Za pravilno kvadriranje binoma $1$ točka, za pravilno uporabo pravila razlike kvadratov $1$ točka, pravilno množenje $1$ točka, pravilno izpostavljanje $1$ točka, pravilno uporabljena Vietova formula $1$ točka.
    \end{solution}
    

    % \newpage
    % 3. vprašanje
    \question
    Faktorizirajte izraze, kolikor je mogoče.
    \begin{parts}
        \part[2]
        $2x^2-30x+108$

        \begin{solution}[4cm]
            $$2x^2-30x+108=2(x-6)(x-9)$$
            Za uporabo Vietovega pravila $1$ točka, pravilen rezultat $1$ točka.
        \end{solution}


        \part[2]
        $f^4-9g^2$

        \begin{solution}[4cm]
            $$f^4-9g^2=(f^2-3g)(f^2+3g)$$
            Za pravilno izpostavljanje $1$ točka, pravilno uporabljena formula razlike kvadratov $1$ točka.
        \end{solution}


        \part[3]
        $c^3-4c^2-c+4$

        \begin{solution}[4cm]
            $$c^3-4c^2-c+4=(c-1)(c+1)(c-4)$$
            Za uporabo pravila faktorizacije štiričlenika 2+2 $1$ točka, pravilna faktorizacija štiričlenika $1$ točka, pravilno razstavljena razlika kvadratov $1$ točka.
        \end{solution}


        \part[3]
        $m^5+27g^6m^2$

        \begin{solution}[4cm]
            $$m^5+27g^6m^2=m^2(m+3g^2)(m^2-6g^2+9g^4)$$
            Za pravilno izpostavljanje skupnega faktorja $1$ točka, uporaba formule razlike kubov $1$ točka, pravilen končen rezultat $1$ točka.
        \end{solution}


        \part[3]
        $81-9b^2+12bc-4c^2$

        \begin{solution}[3cm]
            $$81-9b^2+12bc-4c^2=(9-3b+2c)(9+3b-2c)$$
            Za uporabo pravila faktorizacije štiričlenika 3+1 $1$ točka, pravilna uporaba formule $1$ točka, pravilno razstavljanje razlike kvadratov $1$ točka.
        \end{solution}


    \end{parts}


    %  \newpage

    % 4. vprašanje
    \question
    Določite neznani števki, da bo število $\overline{77777b77777a}$ deljivo
    \begin{parts}
        \part[1]
        s številom $2$.

        \begin{solution}[2cm]
            Za navedbo kriterija za deljivost z $2$ $1$ točka, uporaba kriterija in izračun $a\in\{0,2,4,6,8\}$ $1$ točka.
        \end{solution}


        \part[3]
        s številom $12$.

        \begin{solution}[4cm]
            Pravilno uporabljen kriterij za deljivost s $3$ $1$ točka, pravilno uporabljen kriterij za deljivost s $4$ $1$ točka, 
            pravilen zaključek in zapis kombinacij $(a,b)\in\left\{(2,0),(2,3),(2,6),(2,9),(6,2),(6,5),(6,8)\right\}$ $1$ točka.
        \end{solution}


    \end{parts}



% 5. vprašanje

    \question[4]
    Naj bo $a=3^4\cdot 5^4\cdot 7\cdot 13$, $b=2^2\cdot 3^2\cdot 175$ in $c=36\cdot 625$. 
    Določite največji skupni delitelj $D(a,b,c)$ in najmanjši skupni večkratnik $v(a,b,c)$. Vrednost slednjega lahko pustite zapisano s potencami.

    \begin{solution}[7cm]
        Za zapisa $b=2\cdot 3^3\cdot 5^2\cdot 7$ in $c=2^2\cdot 3^2\cdot 5^4$ po $1$ točka. 
        Izračun $D(a,b,c)=3^2\cdot 5^2=225$ $1$ točka, in $v(a,b,c)=2^3\cdot 3^4\cdot 5^4\cdot 7\cdot 13$ $1$ točka.
    \end{solution}
    
    


    % 6. vprašanje
    \question[3]
    Število $m$ da ostanek $6$ pri deljenju z $9$, število $n$ pa da pri deljenju s $3$ ostanek $2$.
    Pokažite, da je število $3n+m$ deljivo s $3$.
    
    \begin{solution}[3cm]
        Za zapis $m=9k+6$ in $n=3l+2$ $1$ točka, izračun $3n+m=9k+9l+12$ $1$ točka, končen sklep $2a+b=3(3k+3l+4)$ $1$ točka.
    \end{solution}

        




     \newpage

    %dodatno vprašanje
    \bonusquestion[3]
    \textit{Dodatna naloga:} Izračunajte in rezultat zapišite v ustreznem številskem sestavu.
                $$ 3724_{(8)}+ 254_{(6)}= \quad \quad \quad _{(2)} $$

    \begin{solution}[7cm]
        $$ 3724_{(8)}+ 254_{(6)}= 100000111110_{(2)} $$
        Za pravilno zapisan algoritem $1$ točka, zapis števila v sedmiškem sistemu $532_{(10)}=1360_{(7)}$ $1$ točka.
    \end{solution}




\end{questions}



\end{document}